\scanpage{25}\mspage{19}
% @ top, a | text | u,v,w liés
a) $u,v,w$ liés deux à deux

% @ top, b | text | u lié aux éléments
b) $u$ lié aux éléments de $\widetilde P_u$, $v$ et $w$ non liés à ces éléments.

% @ top, c | text | u non lié aux él.
c) $u$ non lié aux él. de \deleted{\unsure{$E$}} $P$. Si $v\in P\wedge_{\mathbb F_2^{P_u}}(\mathbb F_2,\varepsilon)$, il est lié à $x\in P$ sss il est dans l'image de $x$ par $P\longrightarrow P\wedge_{\mathbb F_2^{P_u}}(\mathbb F_2,\varepsilon)$.

% @ upper third, d | text | Deux éléments
d) Deux éléments de $\widetilde P_u$ liés sss ils sont conjugués par $\sigma$, i.e. \deleted{\unsure{qu'ils}} \deleted{\unsure{liés par un}} sont $\in$ des $=$ \unsure{sommands|sommants} $P\wedge_{\mathbb F_2^{P_u}}(\mathbb F_2,p_i)$ $(i\in P_u)$ et sont distincts (i.e. sont distincts et ont $=$ image dans $\widetilde P_u\text{---}P_u$)

% @ upper middle, e | text | Un él. a
e) \deleted{Deux} Un él. $a$ de $\widetilde P_u$ lié à un él. $x$ de $P$ sss, si $a\in P\wedge_{\mathbb F_2^{P_u}}(\mathbb F_2,p_i)$, $a$ est dans l'image de $x$.

% @ upper middle, f | text | x et y
f) $x$ et $y\in P$ sont liés sss, posant $x-y=u$, \deleted{\unsure{$x$}}\inserted{$u$} est de la forme \deleted{\unsure{$e_1,e_2,e$|$e'_1,e_2,\ldots$} $e'_i$ $(i\in P_u)$,}
% @ above the fold, two sums | display | e′ᵢ ; e′₀
\[
e'_i=\sum_{j\in P_u-\{i\}}e_{\unsure{i|j}}\quad\text{\uline{ou}}\quad\deleted{\Sigma}\ e'_0=\sum_{i\in I}e_i
\]

% @ lower half, ruled corollary | text | Corollaire Le groupe des autom.
% @ left of corollary, first vertical | drawing | vertical line beside first paragraph
\noindent\makebox[0pt][r]{\raisebox{-4\baselineskip}[0pt][0pt]{\tikz\draw (0,0)--(0,2.1);}\hspace{.6em}}\uline{Corollaire} Le groupe des autom. \inserted{$W_{\{v,w\}=\alpha}$} de $(C,\{v,w\})$ (où $\{v,w\}$ est une arête marquée) est isom. au groupe des couples $(\rho,\varphi)$, où $\rho$ \inserted{$\in\unsure{\mathfrak S|\mathscr S}_{P_u}$} est un autom. de \deleted{groupe} $P_u$, et $\varphi$ un $\bar\rho$-autom de $P$. Donc on a une suite exacte \inserted{(splittée)}
% @ lower third, exact sequence | diagram | 1, F₂, Wα, S, 1
\diagram{\begin{tikzcd}[ampersand replacement=\&,column sep=large]
1\arrow[r,dash]\&(\mathbb F_2^{P_u})^{\deleted{\unsure{P|F}}}\arrow[r]\&W_\alpha\arrow[r,dash]\&\mathfrak S_{P_u}\arrow[r,dash]\&1
\end{tikzcd}}
% @ lower third, identification | text | et Wα est isom.
% @ left of identification | drawing | second vertical line
\noindent\makebox[0pt][r]{\raisebox{-1.3\baselineskip}[0pt][0pt]{\tikz\draw (0,0)--(0,1.1);}\hspace{.6em}}et $W_\alpha$ est isom. au groupe de Weyl $W_{B_4}$ de $B_4$ $(\simeq\mathrm{so}(9))$

% @ foot, root count and change of situation | text | Ce système de racines
Ce système de racines $B_4$ a 32 racines \inserted{(mais pas ttes de $=$ \unsure{longr.|longueur} !)} qu'on aimerait retrouver plongé dans $E_6$. Il vaut mieux, pour éliminer les longueurs inégales, épingler légèrement la situation,
% archive sheet number: 24
